\section{High-Advantage Responses from the Initial Student}
\label{app:high_advantage_examples}

We present three responses sampled from Qwen3-1.7B-Base before any OPD
update and scored by the Qwen3-4B teacher. All three belong to the
highest-advantage group, A10, in the analysis of
Section~\ref{sec:teacher_selection} and illustrate excessive length and
mechanical repetition already present under $S_{\mathrm{INIT}}$.
All three are incorrect and contain no parseable final answer.

\paragraph{Reading the complete responses.}
Each response consists almost entirely of a repeated string. We specify its
complete content using exact repetition counts: $+$ denotes concatenation
without an added separator, and $\operatorname{repeat}(s,n)$ denotes $n$
consecutive copies of $s$. Quotation marks delimit strings and are not part
of the response; \texttt{\textbackslash n} denotes a newline. Unicode
characters are written as \texttt{\textbackslash uXXXX} escapes. No text is
omitted from these representations. Prompt mathematics below is
typeset for readability. Every prompt ends with: ``Please reason step by
step, and put your final answer within \texttt{\textbackslash boxed\{\}}.''

\subsection{Length-Cap Repetition}

\paragraph{Prompt (AIME24, \texttt{aime24\_001}).}
Every morning Aya goes for a $9$-kilometer-long walk and stops at a coffee
shop afterwards. When she walks at a constant speed of $s$ kilometers per
hour, the walk takes her 4 hours, including $t$ minutes spent in the coffee
shop. When she walks $s+2$ kilometers per hour, the walk takes her 2 hours
and 24 minutes, including $t$ minutes spent in the coffee shop. Suppose Aya
walks at $s+\frac{1}{2}$ kilometers per hour. Find the number of minutes the
walk takes her, including the $t$ minutes spent in the coffee shop.

\paragraph{Response (sample 65).}
\begin{quote}
\small\ttfamily
"\textbackslash uFFFD\textbackslash uFFFD\textbackslash u6001\ T 10T 50 10 1 10 10"\\
\normalfont $+\;\operatorname{repeat}($\texttt{"T 10"}$,\;2043)$
\end{quote}
The response mechanically repeats \texttt{T 10} until it reaches the
8,192-token generation cap. The two initial
\texttt{\textbackslash uFFFD} characters are present in the stored response.

\subsection{Severe Repetition without Reaching the Length Cap}

\paragraph{Prompt (AIME25, \texttt{aime25\_028}).}
Let the sequence of rationals $x_1,x_2,\ldots$ be defined such that
$x_1=\frac{25}{11}$ and
$x_{k+1}=\frac{1}{3}\left(x_k+\frac{1}{x_k}-1\right)$.
$x_{2025}$ can be expressed as $\frac{m}{n}$ for relatively prime positive
integers $m$ and $n$. Find the remainder when $m+n$ is divided by 1000.

\paragraph{Response (sample 110).}
\begin{quote}
\small
$\operatorname{repeat}($\texttt{".slfwith\textbackslash n"}$,\;2154)
+$ \texttt{".slfwith"}
\end{quote}
The response consists of repeated \texttt{.slfwith} lines and ends after
6,465 recorded tokens, showing that severe repetition can occur without
reaching the generation cap.

\subsection{Severe Repetition at the Length Cap}

\paragraph{Prompt (AIME25, \texttt{aime25\_013}).}
Alex divides a disk into four quadrants with two perpendicular diameters
intersecting at the center of the disk. He draws 25 more line segments
through the disk, drawing each segment by selecting two points at random
on the perimeter of the disk in different quadrants and connecting those
two points. Find the expected number of regions into which these 27 line
segments divide the disk.

\paragraph{Response (sample 58).}
\begin{quote}
\small\ttfamily
"\textbackslash u6CD5\ newVal89\textbackslash u6CD5\ newVal8\textbackslash u6CD5\ newVal8"\\
\normalfont $+\;\operatorname{repeat}($\texttt{"\textbackslash u6CD5\ newVal"}$,\;4091)$
\end{quote}
The response repeats the same fragment until it reaches the 8,192-token
generation cap, exhibiting both excessive length and severe repetition.

Together, these cases illustrate how excessively long and repetitive responses can receive
high relative teacher advantage already under $S_{\mathrm{INIT}}$.
